% ================================================
% 文件: 06_导航接收机.tex
% 章节: 第11章 导航接收机
% 所属部分: 接收机技术
% 版本: v1.0
% ================================================
% 本章目录结构（树状）:
% \chapter{导航接收机}
% ├── \section{导航接收机概述}
% ├── \section{GPS接收机原理}
% ├── \section{北斗导航接收机}
% ├── \section{导航接收机设计要点}
% └── \section{导航接收机应用}
% ================================================

\chapter{导航接收机}
\label{chap:navigation_receiver}

\section{导航接收机概述}
\label{sec:nav_receiver_intro}

导航接收机是一种能够接收卫星导航信号并计算位置、
速度和时间信息的设备。
全球卫星导航系统（GNSS）包括GPS、
北斗、
GLONASS、
Galileo等。

导航接收机的主要功能：
\begin{itemize}
    \item \textbf{信号捕获}：搜索并锁定卫星信号
    \item \textbf{信号跟踪}：持续跟踪卫星信号
    \item \textbf{位置计算}：根据卫星信号计算用户位置
    \item \textbf{时间同步}：提供精确的时间信息
\end{itemize}

卫星导航接收机（常称 GNSS 接收机）与前面几章的接收机有本质区别：它接收的并非“强而可懂”的语音或图像，而是淹没在噪声基底之下、功率极弱（典型到达电平约 $-130$ dBm 甚至更低）的扩频信号。其定位思想是基于“已知卫星位置 + 测量信号传播时间”的几何交汇。一颗卫星播发带精确时间戳的测距码与导航电文，接收机比较本地复现码与收到的码之间的时延，即可算得“伪距”；同时利用载波多普勒测速、利用载波相位提高测距精度。

现代导航接收机多为多系统兼容机，
可同时处理 GPS（美国）、
BDS（中国）、
GLONASS（俄罗斯）、
Galileo（欧盟）等系统的信号，
因而可见卫星更多、
定位更可靠、
收敛更快。
其内部通常由天线与前置滤波、
射频下变频、
相关器/基带处理、
定位解算与用户接口几部分组成。
由于卫星信号穿过多径、
电离层与对流层，
并受接收机钟差影响，
接收机必须完成捕获、
跟踪、
伪距/载波观测提取、
误差修正与最小二乘（或卡尔曼滤波）解算等一系列复杂处理，
这正是后续各节的内容。

一套 GNSS 由空间段、
控制段与用户段组成：
空间段是绕地球运行、
不断播发测距码与导航电文的卫星星座；
控制段由地面主控站与监测站组成，
负责卫星钟差、
轨道与完好性的测定与注入；
用户段即各类接收机。
为在任何时刻、
任何地点能看到足够卫星（通常需 4 颗以上用于三维定位加钟差），
中轨（MEO）星座需二十余至三十余颗卫星按多轨道面布设，
低轨增强与地球同步轨道（GEO）卫星则用于增强与完好性广播。

接收机本身按能力分多类：
单频机（仅 L1/B1）成本低、
多用于消费电子；
双频/多频机利用频间组合削弱电离层延迟，
精度与可靠性更高，
多用于测绘与专业领域；
按用途又有导航型、
授时型（强调时间精度）与测姿/定向型（多天线测基线姿态）。
实际使用的最大挑战来自信号极弱与传播环境：
城市峡谷中卫星被遮挡、
信号经建筑反射形成多径，
使伪距含误，
接收机须借助多系统、
惯性融合与多径抑制来兜底。
理解这三段结构与挑战，
是后续原理与设计讨论的背景。

定位解算在几何上等价于“以卫星为球心、
伪距为半径的多球交汇”。
当可见卫星多于四颗时，
接收机用加权最小二乘（或扩展卡尔曼滤波）求解用户位置、
速度、
钟差，
并给出精度因子（DOP）——DOP 越小几何构型越好、
定位越可靠。
城市峡谷、
峡谷或遮挡会使可见卫星骤减、
几何变差，
定位误差增大甚至失效，
此时常借助惯性导航（INS）、
蜂窝与地图匹配做松/紧组合以保连续。

多径是民用接收机的头号误差源之一：
卫星信号经建筑、
地面反射后与直达信号叠加，
使相关峰畸变、
伪距偏移。
抑制手段包括窄相关间距、
多相关器鉴路、
接收机天线扼流圈与抑制仰角低的信号。
高动态载体（飞机、
导弹）还要求载波环在大多普勒变化下不丢锁，
常采用 FLL 辅助 PLL 的跟踪策略。
综上，
导航接收机的解算质量不仅取决于“算得准”，
更取决于“测得稳、
跟得住”。

\section{GPS接收机原理}
\label{sec:gps_principle}

GPS（全球定位系统）是美国开发的卫星导航系统，
由24颗卫星组成。

GPS定位原理：
\begin{itemize}
    \item \textbf{伪距测量}：测量信号传播时间，计算距离
    \item \textbf{三角定位}：至少需要4颗卫星才能确定三维位置
    \item \textbf{误差源}：电离层延迟、对流层延迟、卫星钟差、多路径效应
\end{itemize}

GPS信号结构：
\begin{itemize}
    \item \textbf{C/A码}：粗捕获码，用于快速捕获
    \item \textbf{P码}：精码，用于高精度定位
    \item \textbf{导航电文}：包含卫星轨道信息和时间信息
\end{itemize}

GPS 采用“码分多址”（CDMA）：
每颗卫星使用不同的伪随机噪声（PRN）测距码，
所有卫星在同一频率上发射，
接收机靠相关处理把目标卫星分离出来。
最典型的 C/A 码（粗捕获码）为 1.023Mchip/s 的 Gold 码，
调制在 L1 载频（约 1575.42MHz）上；
P(Y) 码（精码）速率更高、
用于授权用户的精密定位。
导航电文（50bps）播发卫星星历、
历书、
钟差参数与时间，
使接收机能算出每颗卫星在信号发射时刻的精确位置。

伪距测量是定位的基础。设接收机本地钟相对系统时的偏差为 $\delta t$，则第 $i$ 颗卫星的伪距观测为

\[
\rho_i=\sqrt{(x-x_i)^2+(y-y_i)^2+(z-z_i)^2}+c\,\delta t+\varepsilon_i
\]

其中 $(x,y,z)$ 为用户未知位置，$(x_i,y_i,z_i)$ 为卫星位置，$c$ 为光速，$\varepsilon_i$ 包含噪声与各类误差。未知量有 $x,y,z,\delta t$ 共四个，故至少需四颗卫星联立方程才能解出三维位置与钟差；卫星多于四颗时可用最小二乘求得最优解。

由于卫星信号经过电离层（色散延迟）、
对流层、
受卫星与接收机钟差以及地面多径影响，
伪距含系统误差。
工程上通过双频消电离层组合、
导航电文中的修正参数、
差分（DGPS/RTK）等手段削弱之。
载波相位测量精度远高于伪距，
其观测可写为

\[
\phi=\frac{2\pi}{\lambda}\rho+2\pi N+\text{噪声}
\]

其中 $N$ 为整周模糊度（integer ambiguity）；通过求解 $N$ 可实现厘米级 RTK 定位，但整周模糊度的固定是载波相位定位的核心难点。速度则可由载波多普勒频移反推：卫星与接收机的相对径向速度产生 $f_d\approx (v/c)f_0$ 的多普勒，从而估计用户速度。

导航电文中的星历与历书分工明确：
星历是本星短时间内的精确轨道与钟差，
用于精确定位，
需定期更新；
历书是全星座粗略轨道，
用于预测卫星可见性与辅助捕获，
更新较慢。
卫星还播发“健康标志”指示本星是否可用，
接收机据此剔除故障星。
冷启动需先收齐历书再收星历、
耗时较长；
温启动已有历书只需收星历；
热启动连星历与近似位置都在，
最快可数秒定位——这正是 A-GNSS 通过外部下发这些数据加速的原理。

接收机自主完好性监测（RAIM）利用冗余卫星的一致性检验识别“单一卫星故障”：
当某颗卫星观测明显偏离其余卫星的交汇解时，
判定其不可信并剔除，
从而在无外部增强时仍提供一定完好性告警能力。
生命安全应用（航空进近等）对完好性要求极高，
还需星基增强（SBAS）或地基增强（GBAS）提供误差修正与置信边界。
理解电文结构、
启动条件与完好性机制，
是评估导航接收机适用场景的基础。

\section{北斗导航接收机}
\label{sec:beidou_receiver}

北斗卫星导航系统（BDS）是中国自主开发的卫星导航系统。

北斗系统特点：
\begin{itemize}
    \item \textbf{区域覆盖}：北斗二号覆盖亚太地区
    \item \textbf{全球覆盖}：北斗三号覆盖全球
    \item \textbf{特色服务}：短报文通信、精密单点定位
    \item \textbf{兼容性}：与GPS、GLONASS、Galileo兼容
\end{itemize}

北斗信号包括B1I、
B1C、
B2a、
B3I等频段。

北斗系统历经北斗一号（试验）、
北斗二号（区域）到北斗三号（全球）的演进，
目前北斗三号已提供全球定位、
导航与授时（PNT）服务。
其突出特色之一是在导航之外提供无线电测定（RDSS）短报文通信能力，
使用户在无蜂窝网络时也能收发有限报文，
这对远洋、
抢险与野外作业意义重大；
此外还提供精密单点定位（PPP）等增强服务。
北斗采用与 GPS 等兼容的 CDMA 体制与相近频段，
现代多模接收机可同时跟踪北斗与 GPS 等系统，
提升可用性与可靠性。

北斗三号公开服务信号包含 B1I、
B1C、
B2a、
B3I 等多个频点（B1I 约 1561MHz，
B1C 与 GPS L1 同频约 1575.42MHz，
B2a 与 GPS L5 同频约 1176.45MHz，
B3I 约 1268.52MHz，
具体以官方频点为准）。
多频观测使接收机可用频间组合削弱电离层延迟，
提高定位精度。
接收机对北斗信号的捕获、
跟踪与解算流程与 GPS 类似，
差异主要在 PRN 码结构、
导航电文格式与星历模型；
多系统兼容芯片已能统一处理各系统的伪距与载波观测。

为便于工程对照，
下表列出主流 GNSS 系统的典型频点（数值为公开标准值，
具体以各系统接口控制文件为准）。

\begin{table}[htbp]
\centering
\caption{主流 GNSS 典型频点对照}
\begin{tabular}{|p{4.0cm}|p{5.0cm}|p{6.0cm}|}
\hline
\textbf{系统} & \textbf{典型频点} & \textbf{主要公开信号} \\
\hline
GPS（美国） & L1 1575.42MHz & C/A、L2C、L5 \\
\hline
BDS（中国） & B1I 1561.098MHz & B1I、B1C、B2a、B3I \\
\hline
GLONASS（俄） & L1 约 1602MHz（频分） & G1、G2、G3 \\
\hline
Galileo（欧盟） & E1 1575.42MHz & E1、E5a、E5b、E6 \\
\hline
\end{tabular}
\end{table}

北斗三号（BDS-3）在区域服务之外提供全球 RNSS（无线电导航卫星服务），
并保留北斗特色的 RDSS（无线电测定服务）短报文通信。
其公开服务信号含 B1I、
B1C、
B2a、
B3I 等多个频点，
其中 B1C、
B2a 与国际系统（GPS L1/L5、
Galileo E1/E5a）同频或兼容，
便于多系统联合解算与互操作。
BDS-3 还播发 B2b 上的精密单点定位（PPP-B2b）改正信息，
使用户在不依赖地面基准站的情况下获得高精度定位，
是北斗的重要增强特色。

多系统融合是现代接收机的常态：
同时跟踪 GPS、
BDS、
GLONASS、
Galileo 可显著增加可见卫星数，
改善几何精度因子（GDOP），
在城市、
峡谷与高动态场景下提升可用性与连续性。
接收机对北斗信号的处理与 GPS 类似，
差异主要在 PRN 码结构、
导航电文（北斗星历参数、
历书与钟差模型）与频段规划；
多模基带芯片已能统一完成各系统的捕获、
跟踪、
伪距/载波观测提取与联合解算。
随着北斗三号全球组网完成，
兼容北斗已成为各类导航终端的标配能力。

\section{导航接收机设计要点}
\label{sec:nav_receiver_design}

设计导航接收机时需考虑：

\begin{itemize}
    \item \textbf{多频接收}：支持多个频段，提高定位精度
    \item \textbf{低功耗设计}：适用于便携设备
    \item \textbf{抗干扰能力}：抵抗射频干扰
    \item \textbf{快速定位}：缩短首次定位时间（TTFF）
    \item \textbf{集成天线}：小型化、高性能天线
\end{itemize}

\medskip\noindent\textbf{扩频、灵敏度与首次定位时间}

导航接收机的核心难点是“在极弱信号下可靠捕获”。
扩频带来的处理增益是其根本保障：
处理增益等于扩频带宽与信息带宽之比，
用分贝表示为

\[
G_p=10\log_{10}\frac{R_c}{R_d}
\]

其中 $R_c$ 为码片速率，$R_d$ 为信息速率。以 C/A 码（1.023Mchip/s）配 50bps 导航电文为例，处理增益约 43dB，使原本低于热噪声的信号也能被相关积累检出。这就解释了导航接收机为何能在 $-130$ dBm 量级的信号下工作。

灵敏度通常分两级：捕获灵敏度（冷启动能搜到卫星的最低电平，典型约 $-130$ dBm）与跟踪灵敏度（已锁定后能维持的最低电平，可低至 $-160$ dBm 量级，依赖长积分与辅助数据）。首次定位时间（TTFF）则取决于启动条件：热启动（星历与近似位置已知）仅数秒；温启动（有历书）数十秒；冷启动（无任何辅助）可能需数十秒到一分钟内。辅助 GNSS（A-GNSS）通过网络下发星历与粗略时间，可显著缩短 TTFF 并在弱信号下帮助捕获。

\medskip\noindent\textbf{抗干扰与完好性}

导航频带属合法频谱，
但易受窄带压制、
欺骗（spoofing）与多径影响。
接收机采用自适应调零天线、
窄带滤波、
多相关器抗多径以及多系统交叉验证来提升稳健性；
高安全场景还需射频干扰检测与欺骗识别。
完好性（integrity）指系统及时告警定位超差的能力，
民用通过接收机自主完好性监测（RAIM）利用冗余卫星检验一致性，
差分与星基增强（SBAS）则进一步提高可用性与完好性。
低功耗与集成天线则是手持、
穿戴与物联网终端的关键，
常采用单芯片射频基带一体化与低噪放集成设计。

导航接收机的天线与射频前端对性能影响显著。
无源陶瓷贴片天线成本低、
增益有限，
常用于消费电子；
带 LNA 的有源天线与扼流圈天线则用于高精度与测姿，
抑多径、
提增益。
参考时钟决定载波环稳定度：
消费机用 TCXO，
授时与测绘用 OCXO 甚至驯服钟，
频率稳定度直接关系冷启动速度与动态表现。
多频板需多路射频通道与更高算力基带，
成本与功耗随之上升。

抗干扰与防欺骗是专业领域焦点。
压制式干扰（宽带噪声或窄带单音）可用自适应调零天线与窄带滤波抑制；
欺骗（转发或生成假星信号）则更难识别，
需多系统交叉验证、
加密电文认证（如导航电文认证 NMA）与惯性融合检测异常。
接收机还应具备射频干扰（RFI）检测与上报，
便于运维定位干扰源。
对于车规与生命安全应用，
还需冗余天线、
多源融合与功能安全设计，
确保在单源失效时仍可降级运行。

\section{导航接收机应用}
\label{sec:nav_receiver_applications}

导航接收机广泛应用于：

\begin{itemize}
    \item \textbf{交通运输}：汽车导航、船舶导航、航空导航
    \item \textbf{测绘测量}：大地测量、工程测量
    \item \textbf{农业}：精准农业、农机导航
    \item \textbf{物流}：车辆跟踪、货物监控
    \item \textbf{智能手机}：移动定位、地图应用
\end{itemize}

在交通运输领域，
车载与船载导航接收机提供实时位置、
航向与路况引导，
航空则依赖 GNSS 进行航路与进近导航，
并常以惯性导航（INS）组合以弥补卫星信号短时丢失。
测绘与工程测量利用载波相位 RTK 接收机实现厘米级定点，
是城市建设、
形变监测与地籍测量的基础工具。
精准农业中，
农机依据接收机定位自动沿预设航线作业，
实现播种、
施肥、
喷药的变量控制，
节约生产资料。

物流与车队管理依靠接收机回传的位置、
速度与时间戳进行调度与防盗；
智能手机则将 GNSS 与蜂窝、
Wi-Fi、
惯性传感器融合，
提供步行导航、
地图与位置服务。
应急救援、
电力授时、
金融时间戳、
科学观测（如地壳运动、
电离层探测）等同样是导航接收机的重要应用。
随着多系统兼容、
低功耗单芯片与辅助定位的成熟，
导航接收机已从专业设备深入日常生活的方方面面，
成为现代信息社会的基础设施之一。

除位置外，
导航接收机的授时能力同样关键：
以卫星系统时为基准，
可为通信基站、
电力系统、
金融交易与时间戳服务提供纳秒级时间同步，
是现代关键基础设施的“时间锚”。
高精度形变监测利用载波相位 RTK/PPP 长期观测桥梁、
大坝、
滑坡与地壳运动，
能在毫米级发现异常。
自动驾驶与高级驾驶辅助（ADAS）依赖多频 GNSS 与惯性导航（INS）紧组合，
在隧道、
城市峡谷等信号中断时凭惯性维持连续定位。

无人机、
机器人、
精准农业与港口装卸等也把导航接收机作为空间基准；
海洋与航空则依赖其全球覆盖与完好性告警。
可以预见，
随着低轨导航增强、
多频多系统深融合与抗欺骗技术成熟，
导航接收机将在更高精度、
更强完好性与更可靠的抗干扰能力方向上持续演进，
进一步嵌入交通、
能源、
通信与安防等国民经济命脉领域。

授时应用的纵深常被低估：
5G 基站时间同步、
智能电网故障定位、
金融高频交易时间戳、
分布式系统对账，
都依赖 GNSS 提供的统一时标；
一旦授时受扰，
可能引发大面积协同失效，
故关键设施普遍采用多源守时与抗干扰授时。
高精度形变监测已用于高铁、
桥梁、
库坝与滑坡的毫米级长期观测，
成为公共安全的前哨。

低轨（LEO）通信星座的兴起为 PNT 带来新可能：
低轨卫星信号更强、
多普勒更大，
可显著缩短收敛、
增强抗欺骗，
被视为未来弹性 PNT 的重要补充；
同时，
不依赖卫星的自主 PNT（惯性、
天文、
地磁、
蜂窝测距）也在发展，
以应对 GNSS 拒止环境。
导航接收机正从“单一卫星定位”走向“天基+低轨+自主多源”的综合定位授时体系，
支撑自动驾驶、
无人系统与数字社会的时空底座。
